\par\noindent\textbf{(1)} 定义 F(T,H)=U−TS−μ₀HM，微分 dF=−S dT−μ₀M dH，故S=−F\_T、μ₀M=−F\_H。混合偏导交换给 (∂S/∂H)\_T=μ₀(∂M/∂T)\_H；必须在一致的磁能约定下保持μ₀，不可混用 B dM 与 H dB。
\par\noindent\textbf{(2)} 定场热容 C\_H=T(∂S/∂T)\_H。将S视为S(T,H)得 dS=(C\_H/T)dT+μ₀(∂M/∂T)\_H dH，该恒等式直接把难测熵的场导数转化为磁化温度斜率，为磁热实验的基础。
\par\noindent\textbf{(3)} 可逆绝热过程dS=0，于是 (∂T/∂H)\_S=−μ₀T(C\_H)\textasciicircum{}\{-1\}(∂M/∂T)\_H。若通常顺磁相的(∂M/∂T)\_H<0，增加H会升温，退磁dH<0会降温；此方向判断从符号和正热容共同得出。
\par\noindent\textbf{(4)} Curie区M=CH/T，故 M\_T=−CH/T²，绝热磁热系数为 μ₀ C H/(C\_H T)>0。该近似需弱场、足够高温且无相变；近临界或低温自旋饱和时 Curie式失效，还要计入晶格热容和不可逆磁滞。
